\apendice{Plan de Proyecto Software}

\section{Introducción}
El presente anexo recoge los aspectos relacionados con la planificación y viabilidad del proyecto desarrollado. Se describen las principales fases de trabajo realizadas durante el desarrollo del sistema SIEM, así como una estimación de los costes asociados y los aspectos legales considerados.

Aunque el proyecto se ha desarrollado con fines académicos, resulta interesante analizar su viabilidad desde una perspectiva similar a la utilizada en proyectos profesionales de desarrollo software.

\section{Planificación temporal}
El desarrollo del proyecto se ha realizado siguiendo una metodología incremental, avanzando progresivamente desde los componentes más básicos hasta la integración completa del sistema.

Las principales fases del proyecto han sido las siguientes:

\begin{itemize}
\item Análisis inicial de requisitos y definición de objetivos.
\item Diseño de la arquitectura general del sistema.
\item Diseño del entorno simulado del polvorín.
\item Definición del formato de eventos y logs.
\item Desarrollo del módulo de ingesta.
\item Implementación de la base de datos SQLite.
\item Desarrollo del motor de correlación.
\item Implementación del módulo de detección de anomalías.
\item Desarrollo de la API REST mediante FastAPI.
\item Desarrollo del dashboard de monitorización mediante Streamlit.
\item Pruebas, validación y documentación final.
\end{itemize}

La naturaleza incremental del proyecto permitió validar cada módulo de forma independiente antes de proceder a su integración con el resto de componentes.

\section{Estudio de viabilidad}
Antes de iniciar el desarrollo del proyecto se realizó un análisis preliminar para determinar si la solución propuesta era viable desde el punto de vista técnico, económico y legal.

El objetivo de este estudio fue comprobar que las tecnologías seleccionadas permitían desarrollar un sistema SIEM funcional dentro de las limitaciones temporales y de recursos propias de un Trabajo Fin de Grado.

Tras analizar las herramientas disponibles y los requisitos planteados, se concluyó que el proyecto era viable utilizando tecnologías de código abierto, una arquitectura ligera y un entorno simulado que permitiera validar el funcionamiento del sistema sin necesidad de disponer de infraestructuras reales.

A continuación se analizan los aspectos económicos y legales más relevantes relacionados con el desarrollo del proyecto.

\subsection{Viabilidad económica}
Desde el punto de vista económico, el proyecto presenta una elevada viabilidad debido al uso de tecnologías de código abierto y herramientas gratuitas.

Las principales tecnologías utilizadas durante el desarrollo han sido Python, SQLite, FastAPI, Streamlit, Git y GitHub, todas ellas disponibles sin costes de licencia.

La estimación económica más relevante corresponde al tiempo de desarrollo invertido. Considerando un escenario profesional y una dedicación aproximada de 300 horas de trabajo a un coste estimado de 20 euros por hora, el coste de desarrollo ascendería aproximadamente a 6.000 euros.

No obstante, al tratarse de un proyecto académico desarrollado con herramientas gratuitas y sin necesidad de infraestructuras específicas, los costes reales han sido mínimos.

\subsection{Viabilidad legal}
El proyecto ha sido desarrollado utilizando exclusivamente tecnologías de código abierto o herramientas con licencias gratuitas para uso académico.

El entorno utilizado corresponde a una simulación ficticia de un polvorín militar y no incorpora información sensible, procedimientos operativos reales ni documentación clasificada.

Asimismo, se han tenido en consideración principios generales relacionados con la monitorización de eventos y la protección de infraestructuras críticas, evitando en todo momento la inclusión de información que pudiera comprometer la seguridad de instalaciones reales.

Por todo ello, no se identifican impedimentos legales para el desarrollo y utilización académica del sistema presentado en este Trabajo Fin de Grado.

